\chapter{Profil Utilisateur}
% =====================================================================

L'\textbf{Utilisateur} est un profil de \emph{consultation simple}. Il consulte les données opérationnelles (fournisseurs, clients, factures, règlements, avances) afin de retrouver une information, mais il ne saisit ni ne modifie aucune donnée, n'accède pas aux rapports et n'accède pas à l'administration. C'est le profil adapté à un agent qui doit simplement rechercher ou vérifier des éléments.

\section{Périmètre du profil}

\renewcommand{\arraystretch}{1.3}
\begin{longtable}{>{\raggedright\arraybackslash}p{5.4cm} >{\raggedright\arraybackslash}p{9.2cm}}
\toprule
\textbf{Module} & \textbf{Droits de l'Utilisateur} \\
\midrule
\endhead
Fournisseurs & Consulter \\
Factures fournisseurs & Consulter \\
Règlements fournisseurs & Consulter \\
Clients & Consulter \\
Factures clients & Consulter \\
Règlements clients & Consulter \\
Avances clients & Consulter \\
Rapports & Non accessibles \\
Plan comptable, Banques & Non accessibles \\
Paramètres, Utilisateurs, Rôles, Journal & Non accessibles \\
\bottomrule
\end{longtable}
\renewcommand{\arraystretch}{1}

\note{Par rapport au Gestionnaire, l'Utilisateur ne dispose pas du menu \menu{Rapports}. Son menu se limite au tableau de bord et aux rubriques de consultation des fournisseurs et des clients.}

\section{Le tableau de bord}

À la connexion, l'Utilisateur arrive sur le \menu{Tableau de bord}, qui offre une vue d'ensemble~: les quatre indicateurs clés (chiffre d'affaires du mois, factures en attente, dettes fournisseurs, créances clients), le tableau des dernières factures fournisseurs, la situation des banques et les deux graphiques.

\screenshot{images/dashboard_accueil.png}{Tableau de bord de l'Utilisateur.}

\section{Consulter les fournisseurs et leurs opérations}

\begin{itemize}
  \item \menu{Factures Fournisseurs > Fournisseurs} — la liste des fournisseurs, avec recherche, filtre par compte comptable et accès à la fiche détaillée (\bouton{Voir}).
  \item \menu{Factures Fournisseurs > Factures} — la liste des factures d'achat avec leurs montants et statuts, les filtres (fournisseur, statut, période) et le détail d'une facture (\bouton{Voir}).
  \item \menu{Factures Fournisseurs > Règlements} — les règlements regroupés par facture, avec leur détail.
\end{itemize}

\screenshot{images/gestion_fournisseur_liste.png}{Consultation des fournisseurs, factures et règlements par l'Utilisateur.}

\section{Consulter les clients et leurs opérations}

\begin{itemize}
  \item \menu{Factures Clients > Clients} — la liste des clients, avec recherche et accès à leurs factures.
  \item \menu{Factures Clients > Factures} — la liste des factures de prestations, avec filtres et détail d'une facture.
  \item \menu{Factures Clients > Règlements} — les règlements clients regroupés par facture.
  \item \menu{Factures Clients > Avances} — les chèques d'avance reçus et leur consommation.
\end{itemize}

\screenshot{images/factures_clients.png}{Consultation des clients, factures, règlements et avances par l'Utilisateur.}

\section{Gérer son compte}

L'Utilisateur accède à \menu{Mon Profil} (informations personnelles, mot de passe, signature) et peut se déconnecter manuellement~; la déconnexion automatique pour inactivité s'applique (voir le chapitre \og Prise en main \fg{}).

% =====================================================================
